\section{区块保存系统}
\label{sec:chunk_storage}

\subsection{背景}

Minecraft的世界数据持久化需要处理海量方块数据。一个标准世界包含数百万个区块，每个区块存储$16\times16\times256=65536$个方块状态、生物群系和光照数据。存储系统需要在以下需求之间取得平衡：

\begin{enumerate}
    \item \textbf{随机读写性能}：玩家可能在任意位置放置或破坏方块
    \item \textbf{空间效率}：世界数据可能达到数GB甚至数TB
    \item \textbf{数据一致性}：崩溃恢复不应丢失用户进度
    \item \textbf{并发访问}：多线程同时读写不同区域
\end{enumerate}

\subsection{原版局限性}

Java版Minecraft使用Anvil文件格式存储区块数据，该格式基于区域文件（Region File）组织：

\textbf{Anvil文件结构}：每个区域文件覆盖$32\times32$个区块（$512\times512$方块区域）。文件开头是1024个区块偏移量表（每个4字节）， followed by 1024个时间戳表。区块数据以NBT格式序列化后使用Zlib压缩存储。

主要问题包括：

\begin{enumerate}
    \item \textbf{存储粒度过大}：整个区块（$16\times16\times256$）作为一个单元存储，更新单个方块需要重写整个区块数据，产生大量不必要的IO
    \item \textbf{压缩算法效率}：Zlib压缩/解压缩开销较大，区块加载时CPU密集
    \item \textbf{文件碎片化}：区域文件随时间增长会出现内部碎片，已删除区块的空间无法有效回收
    \item \textbf{并发限制}：文件锁机制限制了多线程并发读写能力
    \item \textbf{崩溃恢复风险}：如果写入过程中崩溃，可能丢失整个区块甚至整个区域文件的数据
\end{enumerate}

区块加载时间可分解为：
\begin{equation}
T_{read} = T_{seek} + T_{decompress} + T_{parse}
\end{equation}

其中$T_{seek}$为文件定位时间（区域文件头查找+数据偏移定位），$T_{decompress}$为Zlib解压时间，$T_{parse}$为NBT解析时间。这三部分开销都相当可观。

\subsection{本项目实现}

本项目采用RocksDB作为底层存储引擎，结合Section粒度存储和多层缓存，实现高性能区块持久化。

\subsubsection{RocksDB存储引擎}

RocksDB是Facebook基于LevelDB开发的高性能嵌入式键值存储，使用LSM树（Log-Structured Merge Tree）作为核心数据结构。其优势包括：

\begin{itemize}
    \item \textbf{写入优化}：写入操作先写入MemTable（内存中的跳表），达到阈值后刷入SST文件，写入性能接近内存速度
    \item \textbf{内置压缩}：支持多种压缩算法（Snappy、ZSTD、LZ4等），本项目选用ZSTD，压缩率与Zlib相当但速度快约3倍
    \item \textbf{多线程压缩}：后台压缩和刷盘操作使用独立线程池，不阻塞读写请求
    \item \textbf{原子写入}：WriteBatch支持批量操作的原子性提交
    \item \textbf{快照支持}：提供一致性快照，用于备份和读取一致性
\end{itemize}

关键配置参数：
\begin{itemize}
    \item 块缓存（Block Cache）：256MB，缓存解压后的SST数据块
    \item 行缓存（Row Cache）：64MB，缓存键值对
    \item MemTable大小：64MB，写缓冲区
    \item MemTable数量：4（1个活跃+3个不可变），允许后台刷盘时继续写入
    \item LSM树层级：7层，平衡写入放大和读取性能
    \item Bloom过滤器：每key 10 bit，减少不必要的磁盘读取
\end{itemize}

\subsubsection{Section粒度存储}

本项目将区块划分为16个Section（每个$16\times16\times16$方块），实现细粒度存储：

\textbf{键编码}：SectionKey编码为13字节，包含维度ID（2字节）、区块X坐标（4字节）、区块Z坐标（4字节）、Section Y坐标（1字节）和保留字段（2字节）。大端序确保键的字典序与空间位置顺序一致，有利于范围查询。

\textbf{数据格式}：SectionData序列化为二进制格式，包含：
\begin{itemize}
    \item \textbf{Header}（12字节）：格式版本、标志位、非空方块计数、内容哈希
    \item \textbf{BlockStates}（变长）：4096个u32方块状态ID，ZSTD压缩存储
    \item \textbf{Biomes}（128字节）：64个u16生物群系ID，4$\times$4$\times$4采样
    \item \textbf{SkyLight}（可选，2048字节）：天空光照Nibble数组
    \item \textbf{BlockLight}（可选，2048字节）：方块光照Nibble数组
\end{itemize}

标志位设计允许优化存储：空Section（全是空气）仅存储Header，不存储BlockStates数据；光照数据仅存储非默认值（天空光默认15，方块光默认0）。

\subsubsection{LRU缓存层}

在RocksDB之上实现Section缓存，减少磁盘IO：

缓存容量默认1024个Section，使用\texttt{std::list}维护LRU顺序，\texttt{std::unordered\_map}提供O(1)查找。每个缓存条目包含：
\begin{itemize}
    \item Section数据指针
    \item 脏标记：标识是否需要写回磁盘
    \item 最后访问时间：用于统计和调试
\end{itemize}

缓存命中率统计用于性能分析：
\begin{equation}
HitRate = \frac{Hits}{Hits + Misses}
\end{equation}

脏Section追踪机制支持增量保存：只有被修改的Section才会触发磁盘写入，避免不必要的序列化和压缩开销。

\subsubsection{列族设计}

RocksDB支持列族（Column Family）隔离不同类型的数据，本项目按维度和数据类型划分：

\begin{itemize}
    \item \texttt{sections\_overworld}：主世界Section数据
    \item \texttt{sections\_nether}：下界Section数据
    \item \texttt{sections\_the\_end}：末地Section数据
    \item \texttt{entities\_*}：实体数据（按维度）
    \item \texttt{poi\_*}：POI（兴趣点）数据（按维度）
    \item \texttt{players}：玩家数据
    \item \texttt{snapshots}：快照元数据
\end{itemize}

列族隔离的优势：
\begin{itemize}
    \item 独立压缩配置：Section数据压缩率高，实体数据压缩率低
    \item 范围查询高效：可以单独遍历某个维度的数据
    \item 冷热数据分离：主世界访问频繁，末地访问稀少
\end{itemize}

\subsubsection{异步IO与自动保存}

上层通过\texttt{std::future}异步接口访问存储系统：

\texttt{loadSectionAsync}使用\texttt{std::launch::async}启动异步任务，返回\texttt{future<Result<SectionData*>>}。调用方可以在等待结果的同时执行其他工作。

自动保存机制结合定时触发和阈值触发：
\begin{itemize}
    \item 定时触发：每5分钟检查一次是否需要保存
    \item 阈值触发：脏Section数量达到1000时立即保存
    \item 批量写入：每次保存最多100个Section，避免一次性写入过多导致卡顿
\end{itemize}

\subsection{解决与成果}

\begin{table}[h]
\centering
\caption{存储系统对比}
\label{tab:storage_comparison}
\begin{tabular}{lcc}
\toprule
\textbf{指标} & \textbf{Java Anvil} & \textbf{本项目RocksDB} \\
\midrule
存储粒度 & Chunk (256KB) & Section (16KB) \\
压缩算法 & Zlib & ZSTD \\
索引结构 & 区域文件头 & LSM树 \\
增量更新 & 不支持 & 支持 \\
并发访问 & 文件锁限制 & 多线程原生支持 \\
崩溃恢复 & 区块丢失风险 & WAL保证原子性 \\
\bottomrule
\end{tabular}
\end{table}

Section粒度存储使得单个方块更新的写入量从256KB降至约16KB（减少94\%），配合ZSTD压缩（比Zlib快约3倍），区块加载性能提升约5倍。LSM树的写入优化使得大量区块保存时的吞吐量显著提升，WAL机制保证了崩溃恢复的数据完整性。
